% Original applied explanation; see Chapter04/SOURCES.md.
\subsubsection{ট্যাগ, অ্যাট্রিবিউট ও ভেতরের লেখার ভূমিকা}
\texttt{<a href="notice.html">Notice</a>} অংশে \texttt{a} এলিমেন্ট লিংক বোঝায়, \texttt{href} তার গন্তব্য জানায় এবং Notice হলো ব্যবহারকারী যে লেখা দেখেন। পুরো অংশটি একটি এলিমেন্ট; অ্যাট্রিবিউটটি শুরুর ট্যাগে থাকে।

নেস্টিং পড়বে বাইরের অংশ থেকে ভেতরে। অনুচ্ছেদের ভেতরে গুরুত্ব বোঝানোর অংশ খুললে আগে সেই ভেতরের অংশ, তারপর অনুচ্ছেদ বন্ধ করবে। ব্রাউজার কিছু ভুল নিজে সংশোধন করে ফল দেখাতে পারে; ফল দেখা গেলেই মূল কোডের গঠন সঠিক, এমন সিদ্ধান্ত নেবে না।
